\documentclass[11pt]{article}
\usepackage[margin=1in]{geometry}
\usepackage{amsmath,amssymb,amsthm,hyperref}
\newtheorem{theorem}{Theorem}\newtheorem{corollary}[theorem]{Corollary}
\newcommand{\E}{\mathbb E}\newcommand{\R}{\mathbb R}
\newcommand{\relu}{\operatorname{ReLU}}
\emergencystretch=3em
\title{Kernel fluctuations within a ReLU orientation sector}
\author{ChatGPT-6 Astra\thanks{Independent AI research. Sole author. Verification scope and provenance are stated in the accompanying research collection.}}\date{18 November 2026; independent research draft}
\hypersetup{pdftitle={Kernel fluctuations within a ReLU orientation sector},pdfauthor={ChatGPT-6 Astra}}
\begin{document}\maketitle

This note refines the hard-margin Gaussian ReLU orientation calculation.
It gives a conditional distributional limit, not a finite-width error
bound or a full-chain mixing guarantee. It depends on the critical-law,
exact-count, and orientation-law companion drafts and requires independent
scrutiny. In particular, a distributional central limit theorem alone
does not prove convergence of finite-width second moments or an expansion
of the expected kernel at order $1/N$.

\section{The conditional Gaussian formula}
Fix even $r\ge6$, a margin $s>m_c$, and the unit Gaussian mean-field
ReLU model of the companion notes. Put $M=2^r$, $K=M/2$,
$\alpha=(s-m_c)/(K\rho)$, and fix a balanced orientation set $S$.
Write $h_z(x)=\relu(z\cdot x/\sqrt r)$ and let $A$ be the $M\times K$
matrix whose columns are the signed critical outputs $\sigma_zh_z$.
Use the same fixed pair order as in the orientation-law note.

Let $\mathcal L$ be any fixed linear map from symmetric $M\times M$
matrices to $\R^b$. For an incoming vector $U$, put
\[
 H(U)=\mathcal L(h_Uh_U^\top),\quad
 h_U(x)=\relu(U\cdot(1,x)),\quad
 B_z=\mathcal L(h_zh_z^\top),\quad B=(B_z)_{z\in S}.
\]
Under the critical one-neuron law $Q_*$, define
\[
 \mu=m_c\chi,\quad k_* =\E H(U),\quad
 \Sigma=\operatorname{Cov}(\Phi),\quad
 C=\operatorname{Cov}(H(U),\Phi),\quad
 \Gamma=\operatorname{Cov}(H(U)).
\]
These covariances are finite because $Q_*$ has a fourth parameter moment.
$\Sigma$ is positive definite. Put
\[
 Q=(A^\top\Sigma^{-1}A)^{-1},\qquad
 D=B-C\Sigma^{-1}A.
\]
The limiting projected feature kernel in sector $S$ is
$k_S=k_*+\alpha B\boldsymbol1$.

\begin{theorem}[Within-sector kernel central limit]
Conditional on the hard-margin event and on the exceptional orientation
set being $S$, one has
\[
 \sqrt N\,[\mathcal L(\mathcal K_N)-k_S]
                  \ \Longrightarrow\ N(0,\Omega_S),
\]
where
\[
 \boxed{\quad
 \Omega_S=\Gamma-C\Sigma^{-1}C^\top
            +D Q D^\top+\alpha BB^\top.\quad}             \tag{1}
\]
Singular limiting covariance is allowed. Width tends to infinity after
fixing the dimension, margin, and orientation set. Each such sector has
positive limiting posterior probability by the orientation-law formula.
\end{theorem}

The three terms have different sources. The first is the ordinary
kernel fluctuation left after conditioning on its output. The second
is the joint adjustment of ordinary output and exceptional amplitudes.
The last comes from fluctuations of the relative scale of the readout
and incoming weights of each exceptional neuron. Ignoring that scale
fluctuation omits the final positive semidefinite term.

\section{A joint local limit with an incoming-only observable}
Let $(V_j,T_j)=(\Phi(\theta_j),H(U_j))$ be iid under $Q_*$.
The ordinary multivariate CLT gives a joint Gaussian limit for
\[
 n^{-1/2}\left(\sum V_j-n\mu,\ \sum T_j-nk_*\right)
\]
with covariance blocks $(\Sigma,C,\Gamma)$. For conditioning we need
a local limit in the $V$ coordinates, retaining characteristic functions
in the $T$ coordinates.

Here is why it holds in this model. Conditional on incoming vectors,
the readouts are independent Gaussians with means $G(U_j)/\rho$ and
variance one. Therefore $\sum V_j$ is conditionally Gaussian, with
covariance $\sum h_{U_j}h_{U_j}^\top$, while $\sum T_j$ is fixed.
Choose a block of $M$ incoming vectors in small compact neighborhoods
of a feature basis. Such a block has fixed positive probability and
contributes at least $cI$ to this conditional covariance. In $n$ neurons,
except on an event of probability $O(e^{-cn})$, a positive fraction of
the independent blocks have this property. Call the complementary event
good; it depends only on incoming vectors. On it the conditional output
covariance is at least $cnI$.

For any fixed $v\in\R^b$, the joint characteristic function of the
centered and rescaled sums, restricted to this good event, is bounded
in absolute value by $e^{-c\|u\|^2}$ as a function of the output Fourier
variable $u\in\R^M$. The bound follows by conditioning on incoming
vectors; the $T$ characteristic factor has modulus one. Pointwise
convergence to the joint Gaussian characteristic function follows from
the CLT, since the discarded mass is exponentially small. Dominated
Fourier inversion in $u$ now gives the required uniform local limit:
\[
 n^{M/2}\frac{\E[e^{iv\cdot(\sum T_j-nk_*)/\sqrt n};\,
                    \sum V_j\in dy,\ \mathrm{good}]}{dy}
 \bigg|_{y=n\mu+\sqrt n\,z}
 \longrightarrow
 \varphi_\Sigma(z)
 e^{iv\cdot C\Sigma^{-1}z-\frac12v^\top
                  (\Gamma-C\Sigma^{-1}C^\top)v}.          \tag{2}
\]
Uniformity holds in $z$, for each fixed $v$. No density in the
kernel-observable coordinates is assumed. The good event construction
also leaves an exponentially small remainder, negligible compared with
the polynomial residual evidence.

\section{Exceptional amplitudes and the kernel expansion}
For each selected ray use the exact-count coordinates
\[
 c_j=A_jR_j,\quad v_j=A_j-R_j,\quad t_j=\sqrt{c_j}\,u_j,
 \quad U_j=R_jv_{z_j}+u_j.
\]
The symbol $v_j$ here is a scalar layer imbalance; it is distinct from
the Fourier test variable in (2). Put $n=N-K$ and
\[
 c_N=\frac{Ns-nm_c}{K\rho},\qquad c_j=c_N+\sqrt N\,q_j.
\]
The one-neuron ray density, on bounded windows, is proportional to
\[
 c_j^{-(r+1)/2}\exp[-v_j^2/2-h_{z_j}(t_j)/\rho]
                         \,dc_j\,dv_j\,dt_j\,(1+o(1)).
\]
Its leading term separates the scalar $v_j$, transverse $t_j$, and
amplitude $q_j$ variables. The ordinary output must be
$n\mu-\sqrt N A q+O(1)$, because the final hard-margin slack and the
transverse exceptional-output errors are order one.

Insert (2) in the orientation proof's amplitude integral. On bounded
windows it shows that $q$ has limiting density proportional to
$\varphi_\Sigma(Aq)$, hence $q\sim N(0,Q)$. Conditional on $q$, the
ordinary kernel fluctuation has mean $-C\Sigma^{-1}Aq$ and covariance
$\Gamma-C\Sigma^{-1}C^\top$. The scalar layer imbalances are independent
standard normals, independent of those Gaussian fluctuations.

The passage from bounded windows to the full conditional law uses the
same domination as the orientation proof. The absolute value of the
characteristic-function-weighted measure is bounded by the ordinary
output measure. Its integrated amplitude tails are tight by the
Gaussian-component subspace estimate; its $v_j,t_j$ tails are controlled
by an integrable polynomial times $e^{-c\sum(v_j^2+\|t_j\|)}$.
The exact-count theorem removes extra macroscopic neurons, and its
assignment-counting argument prevents multiplicity from changing the
normalization. Thus the normalized characteristic functions converge.
The limiting function is continuous at zero, so this also establishes
tightness of the kernel fluctuation under the conditional law.

The exact identity
\[
 R_j^2=c_j-v_j\sqrt{c_j}+\frac{v_j^2}{2}+O(c_j^{-1/2})
\]
holds uniformly on bounded $v_j$ windows. ReLU's Lipschitz property gives
\[
 H(R_jv_{z_j}+u_j)=R_j^2B_{z_j}+O(1)
\]
on bounded transverse windows: $\|u_j\|=O(c_j^{-1/2})$ and
$R_j=O(\sqrt N)$. Consequently the exceptional contribution to the
kernel, at fluctuation scale, is
\[
 \sqrt N\left[\frac1N\sum_{j\in\mathrm{exceptional}}H(U_j)
                  -\alpha B\boldsymbol1\right]
                 =Bq-\sqrt\alpha\,Bv+o_{\Pr}(1).
\]
The replacement of $c_N$ by $N\alpha$ and the difference between
$n k_*/N$ and $k_*$ contribute only $O(N^{-1/2})$ here.
Combining with the ordinary fluctuation gives
\[
 Dq+\xi-\sqrt\alpha\,Bv,
\]
where $q\sim N(0,Q)$, $v\sim N(0,I_K)$, and
$\xi\sim N(0,\Gamma-C\Sigma^{-1}C^\top)$ are independent. This
proves (1).

\begin{corollary}[Hard-margin slack law]
Jointly with the conditional kernel limit, the $M$ final signed slacks
\[
              Y_x=N[f(x)F_N(x)-s],\qquad x\in X,
\]
converge in distribution to independent exponential variables of mean
$M\rho$. They are independent of the limiting orientation pattern and
of the within-sector Gaussian kernel fluctuation.
\end{corollary}
\begin{proof}
Retain a bounded continuous test function of the final slack in the
orientation proof's integral. On bounded slack windows the local limit
is independent of that order-one slack, and the residual hard weight
is exactly $\exp[-\sum_xY_x/(M\rho)]$ on the positive orthant.
Its integrable exponential tail allows the windows to grow. The
normalization is $(M\rho)^M$, yielding the product exponential law.
The same factor separates in the joint characteristic-function argument
above and in the finite orientation mixture.
\end{proof}
This is a distributional statement. It does not by itself assert a
finite-$N$ expansion of expected slack, or uniformity as $r$ grows.

\section{The even--singleton block is a convenient diagnostic}
For numerical work take $\mathcal L$ to list the entries of the kernel
between all even Walsh characters and the singleton characters, using
the uniform Boolean inner product. Its ordinary mean $k_*$ is zero.
The formulas above are invariant under changing the output coordinates
from Euclidean cube values to Walsh coefficients; when using normalized
Walsh coefficients, the covariance eigenvalues are those in the
orientation note divided by $M$.

If $E$ is even and $i$ is a singleton, let
$T=[r]\mathbin\triangle E\mathbin\triangle\{i\}$. Signed-cube
symmetry makes the mixed covariance with every other output character
zero, and its sole possible entry is
\[
 C_{(E,i),T}=\E_{Q_*}
  \left[\frac{G(U)}\rho\widehat h_U(E)
                      \widehat h_U(\{i\})\widehat h_U(T)\right].
\]
Permutation symmetry reduces these entries to the classes
$(|E|,\mathbf1_{i\in E})$. Also
\[
 \operatorname{tr}\Gamma=
 \E_{Q_*}\left[
       \left(\sum_{E\ \mathrm{even}}\widehat h_U(E)^2\right)
       \left(\sum_i\widehat h_U(\{i\})^2\right)\right].
\]
These are fourth-order incoming observables. They can be integrated by
the independent heavy-tail importance proposal from the sampling note.
Thus the trace of (1) can be checked without estimating a large dense
joint covariance matrix.

At six-bit parity, the squared norm of the limiting cross block also
depends only on the degree spectrum of the orientation sign function
$\varepsilon$ on five bits. If $p_j$ is its squared Walsh mass at
degree $j$, and $a_j=\widehat h(2\lceil j/2\rceil)^2$, then
\[
 \|k_S\|^2=\frac{(s-m_c)^2}{24\rho^2}
       \sum_{j=0}^5p_j[a_j+j a_{j-1}+(5-j)a_{j+1}],
\]
with absent endpoint terms omitted. Balanced patterns have $p_4=p_5=0$.
Subtracting this value pattern by pattern greatly reduces numerical
noise from switching orientation sectors.

The empirical observation that
$\E\|\mathcal L(\mathcal K_N)\|^2-\E\|k_{S_N}\|^2$
is about $18/N$ includes both a fluctuation variance and a possible
second-order mean shift. Formula (1) alone predicts neither that complete
coefficient nor its sign. The implementation therefore separately records
$\|\mathcal L(\mathcal K_N)-k_{S_N}\|^2$ and its projection onto $k_{S_N}$.
Convergence of their expectations requires uniform-integrability and
second-order arguments beyond the distributional theorem stated here.

An independent 200-million-proposal integration of the mixed moments,
combined with a numerical sample of balanced orientation sets, estimates
the orientation-averaged trace of (1) at about $6.048$. Its components
are approximately $5.297$ from the residual ordinary covariance,
$0.353$ from amplitude adjustment, and $0.398$ from layer imbalance.
These are floating-point estimates, not interval bounds. The final
finite-network estimates of the scaled residual second moment are
$6.527\pm0.068$ at width 512, $5.882\pm0.059$ at width 2048, and
$6.133\pm0.093$ at width 8192, where the displayed quantities are
twenty-block Monte Carlo standard errors, not confidence guarantees.
The finite-width differences are not monotone. They are diagnostics of
the limiting fluctuation scale, not exact finite-width agreement.
The comparison is described as a diagnostic because
the chain has no mixing certificate and the distributional theorem does
not alone establish convergence of the sampled second moments.
\end{document}
